\section*{Chapitre \ref{ch:g10:signals-and-waves} --- Signaux et ondes}

\begin{solution}{exo:g10:signals-and-waves:1}
$f = 1/T$~: (a) $1/\qty{0.020}{s} = \qty{50}{Hz}$~;
(b) $1/\qty{4.0e-3}{s} = \qty{250}{Hz}$~;
(c) $1/\qty{5.0e-5}{s} = \qty{20}{kHz}$.
\end{solution}

\begin{solution}{exo:g10:signals-and-waves:2}
$T = 1/f$~: (a) \qty{20}{ms}~; (b) $1/440 = \qty{2.3}{ms}$~;
(c) $1/\qty{2.4e9}{Hz} = \qty{0.42}{ns}$.
\end{solution}

\begin{solution}{exo:g10:signals-and-waves:3}
\qty{12}{Hz}~: \omterm{def:g10:signals-and-waves:ultrasound}{infrason}. \qty{440}{Hz} et \qty{18}{kHz}~: audible
(le second seulement pour de jeunes oreilles). \qty{40}{kHz} et
\qty{5}{MHz}~: \omterm{def:g10:signals-and-waves:ultrasound}{ultrason}.
\end{solution}

\begin{solution}{exo:g10:signals-and-waves:4}
$T = 5.0 \times \qty{2}{ms} = \qty{10}{ms}$, donc $f = \qty{100}{Hz}$~;
\omterm{def:g10:signals-and-waves:amplitude}{amplitude} $4.0 \times \qty{0.5}{V} = \qty{2.0}{V}$.
\end{solution}

\begin{solution}{exo:g10:signals-and-waves:5}
$d = 340 \times 6.0 = \qty{2040}{m} \approx \qty{2.0}{km}$. La lumière
couvre cela en $2040/(3.00\times10^{8}) \approx \qty{7}{\micro s}$, un
million de fois moins que \qty{6.0}{s}~: l'éclair marque l'instant du
coup.
\end{solution}

\begin{solution}{exo:g10:signals-and-waves:6}
Le \omterm{def:g10:signals-and-waves:sound}{son} descend \emph{et remonte}, couvrant $2d$~:
$d = \frac{1500 \times 0.60}{2} = \qty{450}{m}$. Au-dessus de la fosse~:
$t = 2d/v = 2 \times 1200/1500 = \qty{1.6}{s}$.
\end{solution}

\begin{solution}{exo:g10:signals-and-waves:7}
$d_1 = \frac{3.00\times10^{8} \times 2.4\times10^{-4}}{2} =
\qty{36.0}{km}$ et $d_2 = \qty{35.7}{km}$~: l'avion a couvert
\qty{300}{m} en \qty{1.0}{s}, donc $v = \qty{300}{m/s} \approx
\qty{1100}{km/h}$, en approche.
\end{solution}

\begin{solution}{exo:g10:signals-and-waves:8}
$T = \qty{20}{mm}/\qty{25}{mm/s} = \qty{0.80}{s}$, donc
$f = \qty{1.25}{Hz}$~: $75$ battements par minute. Pour $100$
battements par minute, $T = \qty{0.60}{s}$, c'est-à-dire
$25 \times 0.60 = \qty{15}{mm}$~: la tachycardie apparaît en dessous
de \qty{15}{mm}.
\end{solution}

\begin{solution}{exo:g10:signals-and-waves:9}
$T = 1/200 = \qty{5.0}{ms}$~; deux cycles durent \qty{10}{ms}, étalés
sur $10$ divisions~: \omterm{def:g10:signals-and-waves:oscilloscope}{base de temps} \qty{1}{ms} par division.
La crête doit tenir en $4$ divisions~: \omterm{def:g10:signals-and-waves:oscilloscope}{gain} d'au moins
$3.0/4 = \qty{0.75}{V}$ par division --- prendre \qty{1}{V} par
division (crête $3{,}0$ divisions en haut).
\end{solution}

\begin{solution}{exo:g10:signals-and-waves:10}
$d = v\,t/2$~: paroi avant
$\frac{1540 \times 4.0\times10^{-5}}{2} = \qty{3.1}{cm}$~; paroi
arrière $\frac{1540 \times 6.0\times10^{-5}}{2} = \qty{4.6}{cm}$~;
épaisseur environ \qty{1.5}{cm}.
\end{solution}

\begin{solution}{exo:g10:signals-and-waves:11}
\omterm{def:g10:signals-and-waves:sound}{Son}~: $30/340 = \qty{0.088}{s}$. Radio~:
$3.0\times10^{6}/(3.00\times10^{8}) = \qty{0.010}{s}$. L'auditeur radio
lointain entend le tambour d'abord, d'environ \qty{78}{ms}.
\end{solution}

\begin{solution}{exo:g10:signals-and-waves:12}
En émettant pendant \qty{3.0}{ms} la chauve-souris est sourde aux
\omterm{rem:g10:signals-and-waves:honest}{échos}, c'est-à-dire à tout ce qui est dans
$d = \frac{340 \times 3.0\times10^{-3}}{2} = \qty{0.51}{m}$. Papillon
à \qty{1.7}{m}~: $t = 2 \times 1.7/340 = \qty{10}{ms}$. En se
rapprochant, l'\omterm{rem:g10:signals-and-waves:honest}{écho} revient de plus en plus tôt~; l'impulsion
doit se terminer avant qu'il n'arrive, donc elle doit se raccourcir.
\end{solution}

\begin{solution}{exo:g10:signals-and-waves:13}
$d = \frac{340 \times 1.5}{2} = \qty{255}{m}$. À \qty{155}{m}~:
$t = 2 \times 155/340 = \qty{0.91}{s}$. \omterm{rem:g10:signals-and-waves:honest}{Écho} et claquement se
confondent lorsque $t < \qty{0.1}{s}$, c'est-à-dire à moins de
$d < \frac{340 \times 0.1}{2} = \qty{17}{m}$.
\end{solution}

\begin{solution}{exo:g10:signals-and-waves:14}
\emph{1.} $T = 3.0 \times \qty{5}{ms} = \qty{15}{ms}$,
$f = 1/0.015 \approx \qty{67}{Hz}$.
\emph{2.} $\qty{15}{ms}/\qty{2}{ms} = 7{,}5$ divisions.
\emph{3.} L'écran montre \qty{50}{ms}, puis \qty{20}{ms}~: $3$ cycles
complets ($50/15 = 3{,}3$), puis $1$ ($20/15 = 1{,}3$).
\end{solution}

\begin{solution}{exo:g10:signals-and-waves:15}
\emph{1.} $t = 2 \times 1.2\times10^{6}/(3.00\times10^{8}) =
\qty{8.0}{ms}$.
\emph{2.} $2 \times 1.2\times10^{6}/(2.0\times10^{8}) = \qty{12}{ms}$.
\emph{3.} La fibre ne court pas en ligne droite entre les deux
machines, et chaque routeur et serveur sur le chemin ajoute un délai
de traitement.
\end{solution}

\begin{solution}{pb:g10:signals-and-waves:1}
\textbf{1.} $d = 340 \times 9.0 = \qty{3060}{m} \approx \qty{3.1}{km}$.

\textbf{2.} $3060/(3.00\times10^{8}) \approx \qty{10}{\micro s}$, un
million de fois plus court que \qty{9.0}{s}~: l'éclair marque l'instant
du coup.

\textbf{3.} En \qty{3}{s} le \omterm{def:g10:signals-and-waves:sound}{son} couvre $340 \times 3 =
\qty{1020}{m} \approx \qty{1}{km}$~: les secondes divisées par trois
donnent les kilomètres.

\textbf{4.} $d = 340 \times 6.0 = \qty{2040}{m}$. L'orage a refermé
\qty{1020}{m} en \qty{300}{s}~: $v = \qty{3.4}{m/s} \approx
\qty{12}{km/h}$, en direction du bateau.

\textbf{5.} L'extrémité proche s'entend après $2000/340 = \qty{5.9}{s}$,
l'extrémité lointaine après $5000/340 = \qty{14.7}{s}$~: le grondement
dure environ \qty{8.8}{s}.

\textbf{6.} $d = 1500 \times 0.90/2 = \qty{675}{m}$.

\textbf{7.} $d = 1500 \times 0.20/2 = \qty{150}{m}$.

\textbf{8.} Deux obstacles~: un banc de poissons à $1500 \times 0.16/2 =
\qty{120}{m}$ et le fond à \qty{150}{m}. Le banc nage \qty{30}{m}
au-dessus du fond.

\textbf{9.} L'\omterm{rem:g10:signals-and-waves:honest}{écho} doit revenir avant le ping suivant~:
$d_{\max} = 1500 \times 0.50/2 = \qty{375}{m}$. Depuis une eau plus
profonde l'\omterm{rem:g10:signals-and-waves:honest}{écho} arrive \emph{après} le ping suivant et lui
est attribué~: l'affichage montre un faux fond beaucoup trop peu
profond.

\textbf{10.} Dans l'air, $340 \times 0.60/2 = \qty{102}{m}$ au lieu de
\qty{450}{m}. Un délai ne devient une distance qu'une fois que l'on
connaît le porteur et sa vitesse dans le milieu traversé.

\textbf{11.} $f = 1/0.86 = \qty{1.16}{Hz}$, c'est-à-dire
$1.16 \times 60 \approx 70$ battements par minute.

\textbf{12.} $25 \times 0.86 = \qty{21.5}{mm}$.

\textbf{13.} $T = 15/25 = \qty{0.60}{s}$~: $100$ battements par minute.

\textbf{14.} $T = 60/50 = \qty{1.2}{s}$~; espacement
$25 \times 1.2 = \qty{30}{mm}$.

\textbf{15.} Non~: la \omterm{def:g10:signals-and-waves:period}{période} dérive d'un battement à
l'autre avec l'effort, la respiration et le stress. Le médecin lit la
\omterm{def:g10:signals-and-waves:period}{période} (le rythme), sa régularité, et la forme de chaque
cycle --- le diagnostic est dans le \omterm{def:g10:signals-and-waves:signal}{signal}.

\textbf{16.} Radio~: $5.0\times10^{4}/(3.00\times10^{8}) =
\qty{0.17}{ms}$ --- instantané pour les sens humains.
\omterm{def:g10:signals-and-waves:sound}{Son}~: $50\,000/340 \approx \qty{147}{s}$, deux minutes et demie.

\textbf{17.} $t = 2.02\times10^{7}/(3.00\times10^{8}) = \qty{67}{ms}$.

\textbf{18.} $c \times \qty{1.0}{\micro s} = \qty{300}{m}$. Pour une
fix à \qty{3.0}{m}~: $3.0/(3.00\times10^{8}) = \qty{10}{ns}$ ---
c'est pourquoi les \omterm{def:g10:universal-gravitation:satellite}{satellites} GPS portent des horloges
atomiques.

\textbf{19.} Les \omterm{def:g10:universal-gravitation:satellite}{satellites} sont dans le vide, et le
\omterm{def:g10:signals-and-waves:sound}{son} a besoin d'un milieu~: aucun \omterm{def:g10:signals-and-waves:signal}{signal}
ne quitterait le \omterm{def:g10:universal-gravitation:satellite}{satellite}. (Même en accordant de
l'air tout le trajet, \qty{20200}{km} à \qty{340}{m/s} font plus de
\qty{16}{h} --- un point de position en retard de plusieurs heures.)

\textbf{20.} \qty{340}{m/s} (\omterm{def:g10:signals-and-waves:sound}{son} dans l'air),
\qty{1500}{m/s} (\omterm{def:g10:signals-and-waves:sound}{son} dans l'eau de mer) et
$c = \qty{3.00e8}{m/s}$ (radio et lumière)~; la loi est $d = v\,t$,
divisée par deux pour un \omterm{rem:g10:signals-and-waves:honest}{écho}. L'habitude~: connaître
son porteur, connaître sa vitesse dans le milieu, chronométrer le
délai --- un délai est une distance déguisée.
\end{solution}
